\begin{abstract}
Tool-using LLM agents are increasingly exposed to indirect prompt injections and fabricated tool outputs that hijack otherwise benign tasks, yet most defenses either guard individual prompts or require retraining the underlying model. We present \method{}, a detector-driven defense pipeline that treats an agent trajectory as a heterogeneous semantic DAG, detects compromised executions with a lightweight graph anomaly detector, localizes the risky execution subgraph, and converts the localized evidence---filtered by an LLM postmortem judge that never observes ground-truth labels---into reusable safety experiences that are injected as advisory memory into a replayed agent. We evaluate on \tool{}, a benchmark of 2{,}000 executable virtual-tool tasks with 2{,}218 sanitized trajectories and aligned semantic graphs (1{,}521 normal, 697 abnormal) under a fixed 7:1.5:1.5 grouped split with a 15\% abnormal ratio on all splits. Under a paired, deterministic recorded-trajectory replay protocol in which every agent step passes an execution controller, injecting detector-derived experiences reduces attack success on the treated high-risk cases from 46.2\% to 0.0\% (exact McNemar $p=0.031$) and lowers panel-level ASR from 26.0\% to 21.0\%, with gate decisions and judge verdicts produced by model families disjoint from the attacked agent. We report detection quality (ROC/PR), defense effectiveness (paired ASR with McNemar tests), ablations separating detector-derived experiences from generic safety knowledge, localization quality, and per-component cost. Our results indicate that graph-level anomaly signals localized to a handful of trajectory nodes are sufficient to synthesize compact advisory experiences that measurably harden agents against trajectory-level attacks without modifying the agent model.
\end{abstract}